\documentclass[11pt]{article}
\usepackage[T1]{fontenc}
\usepackage{amsmath,amssymb,amsthm}
\usepackage[margin=1in]{geometry}
\usepackage[hidelinks]{hyperref}\hypersetup{pdftitle={A Breakpoint-Free Proof of a Quantile-Ladder Identity in Sun's Catalan-Constant Argument},pdfauthor={Xiaojian Zeng},pdfsubject={A structural proof of a generalized-quantile ladder identity},pdfkeywords={generalized inverse, quantile, arithmetic ladder, Catalan constant}}

\newtheorem{theorem}{Theorem}
\newtheorem{lemma}[theorem]{Lemma}

\newcommand{\ind}{\mathbf 1}
\newcommand{\R}{\mathbb R}
\newcommand{\Z}{\mathbb Z}

\title{A Breakpoint-Free Proof of a Quantile-Ladder Identity\\
in Sun's Catalan-Constant Argument}
\author{Xiaojian Zeng\\
PBC School of Finance, Tsinghua University, 100190\\
\texttt{xiaojianzeng@foxmail.com}}
\date{}

\begin{document}
\maketitle

\begin{abstract}
Sun's Lemma~6.6 asserts an identity for the lower-mass cost of a continuum
of arithmetic ladders and describes its proof as an exact finite
verification over a breakpoint partition.  We prove the full identity by a
uniform argument.  The main observation is that an integer-valued base law
whose support has oscillation at most three produces a two-step ladder whose
generalized quantile shifts by two after one unit of mass.  A first-moment
calculation then determines every unit-block cost.  The step functions in
Sun's lemma satisfy the support hypothesis for every parameter value,
including all breakpoints, and have mean \(1-2v\).  This yields the claimed
formula for every \(v\in[0,1)\) and every integer \(n\geq1\), without a cell
enumeration.
\end{abstract}

\section{Statement and context}

For a real number \(x\), write \(\{x\}=x-\lfloor x\rfloor\).  Fix
\(v\in[0,1)\), and set
\[
 a=\{20v\},\qquad b=\{41v\},\qquad
 z(y)=(\tfrac12-y)\bmod 1\quad(0\leq y<1).
\]
Define the integer-valued step function
\[
 h_v(y)=2\lfloor20v\rfloor-\lfloor41v\rfloor
       +2\ind_{\{z(y)<a\}}-\ind_{\{z(y)<v\}}
       +\ind_{\{y>b\}}.                                      \tag{1}
\]
On \(X=[0,1)\times\Z_{\geq0}\), equipped with Lebesgue measure in the
first coordinate and counting measure in the second, let
\[
 c_v(y,r)=h_v(y)+2r.
\]
For \(s\geq0\), let \(K_v(s)\) be the infimum of
\(\int_Ac_v\) over measurable \(A\subseteq X\) of mass \(s\).  Equivalently,
if \(F_v(x)\) is the mass of points with cost at most \(x\) and
\[
 \lambda_v(s)=\inf\{x\in\R:F_v(x)\geq s\}\quad(s>0),             \tag{2}
\]
then \(K_v(s)=\int_0^s\lambda_v(t)\,dt\).  This is the usual generalized
inverse convention~\cite{EmbrechtsHofert2013}; it also handles atoms and
partial selection of a tie.  Values of the indicators at their finitely many
endpoints do not affect the measure or \(K_v\).

Sun introduced (1) in Section~6 of~\cite{Sun2026Catalan} and stated the
following identity as Lemma~6.6:
\[
 K_v(n+v)=n^2+K_v(1+v)-1
 \quad\text{for all \(n\in\Z_{\geq1}\) and \(v\in[0,1)\)}.       \tag{3}
\]
The printed proof sorts finitely many ladders and refers to an exact
verification on the later breakpoint partition; the cell ledger is not part
of the article.  The argument below instead explains the shift uniformly.
The distinction between a deductive structural proof and exact computations
used only as diagnostics is also explicit in a recent reconstruction note of
Zeng~\cite[Section~4]{Zeng2026Multiplicity}; that paper concerns a different
combinatorial problem, and is cited here only for this verification
principle.

\section{The ladder lemma}

If an integer-valued function has essential range contained in
\([L,H]\cap\Z\), call \(H-L\) its support oscillation.

\begin{lemma}[Two-step ladder]
Let \(\eta\) be a probability measure supported on
\([L,H]\cap\Z\), where \(H-L\leq3\), and form the locally finite measure
\[
 \mu=\sum_{r\geq0}T_{2r}\eta,
\]
where \(T_d\) denotes translation by \(d\).  Put
\[
 F(x)=\mu(({-}\infty,x]),\qquad
 \lambda(s)=\inf\{x:F(x)\geq s\},\qquad
 K(s)=\int_0^s\lambda(t)\,dt.
\]
If \(m=\int x\,d\eta(x)\), then
\[
 \lambda(s+1)=\lambda(s)+2\quad(s>1),                           \tag{4}
\]
and
\[
 K(s+1)-K(s)=2s+m\quad(s\geq1).                                \tag{5}
\]
\end{lemma}

\begin{proof}
For \(x\geq H-2\), each progression
\(h,h+2,h+4,\ldots\) contributes exactly one rung, with its base mass,
to \((x,x+2]\).  Consequently
\[
 F(x+2)=F(x)+1.                                                  \tag{6}
\]
The same statement holds with both lower rays open.  Moreover,
\(F(H-2)\leq1\).  Indeed, every rung with index at least one is at least
\(L+2\geq H-1\), so only base rungs can lie at or below \(H-2\).

Let \(s>1\) and \(q=\lambda(s)\).  Then \(q>H-2\), and the two quantile
inequalities are
\[
 \mu(({-}\infty,q))<s\leq F(q).
\]
Applying the open and closed versions of (6) transports these inequalities
to
\[
 \mu(({-}\infty,q+2))<s+1\leq F(q+2).
\]
Thus \(q+2=\lambda(s+1)\).  This proves (4), including the case in which
\(s\) lies at an atom boundary.

It remains to determine the integral over a unit quantile block.  Let \(A\)
and \(B\) be the pushforwards under \(\lambda\) of Lebesgue measure on
\((0,1]\) and \((1,2]\), respectively.  The generalized-inverse equivalence
\[
 \lambda(s)\leq x\quad\Longleftrightarrow\quad s\leq F(x)        \tag{7}
\]
shows that \(\mu\) is the pushforward of Lebesgue measure on
\((0,\infty)\) under \(\lambda\).  Iterating (4) on the unit blocks gives
\[
 \mu=A+\sum_{j\geq0}T_{2j}B.                                   \tag{8}
\]
The original ladder construction also gives \(\mu=\eta+T_2\mu\).
Subtracting the translate of (8) from (8), which is legitimate on bounded
sets because all sums there are finite, yields the finite signed-measure
identity
\[
 \eta=A+B-T_2A.                                                  \tag{9}
\]
Both \(A\) and \(B\) have mass one.  Taking first moments in (9) therefore
gives \(\int x\,dB(x)=m+2\), and hence \(K(2)-K(1)=m+2\).

Write \(D(s)=K(s+1)-K(s)\).  For \(1\leq s\leq2\), cancellation of the
common part of the intervals and (4) give
\[
 D(s)-D(1)=\int_1^s\bigl(\lambda(t+1)-\lambda(t)\bigr)\,dt
            =2(s-1).
\]
Likewise, \(D(s+1)=D(s)+2\) for every \(s\geq1\).  Translating the preceding
formula by integer units proves \(D(s)=2s+m\) for all \(s\geq1\), which is
(5).
\end{proof}

\section{Application to the step function}

Let \(\eta_v\) be the distribution of \(h_v(y)\) for uniform
\(y\in[0,1)\).  The intervals \([0,a)\) and \([0,v)\) are nested.  Hence
\[
 p(z)=2\ind_{\{z<a\}}-\ind_{\{z<v\}}
\]
takes values in \(\{0,1,2\}\) when \(a\geq v\), and in
\(\{-1,0,1\}\) when \(a<v\).  Its oscillation is therefore at most two.
Adding \(\ind_{\{y>b\}}\) increases the oscillation by at most one, while
the floor term in (1) is constant in \(y\).  Thus
\[
 \operatorname{osc}(h_v)\leq3                                  \tag{10}
\]
for every \(v\in[0,1)\).  This includes coincident endpoints and every
floor or fractional-part breakpoint.

The map \(y\mapsto z(y)\) preserves Lebesgue measure.  Write
\[
 q=\lfloor20v\rfloor,\quad a=20v-q,\qquad
 r=\lfloor41v\rfloor,\quad b=41v-r.
\]
Direct integration of (1) gives
\begin{align*}
 m_v:=\int_0^1h_v(y)\,dy
   &=2q-r+2a-v+(1-b)\\
   &=1-2v.                                                       \tag{11}
\end{align*}
The lemma applies to \(\eta_v\).  Taking \(s=j+v\geq1\) in (5), and using
(11), yields for every integer \(j\geq1\)
\[
 K_v(j+1+v)-K_v(j+v)=2(j+v)+1-2v=2j+1.                         \tag{12}
\]
For \(n=1\), (3) is tautological.  For \(n\geq2\), summing (12) from
\(j=1\) to \(j=n-1\) gives
\[
 K_v(n+v)-K_v(1+v)=\sum_{j=1}^{n-1}(2j+1)=n^2-1,
\]
which proves (3) in its full stated range.

\section*{Assistance disclosure}

AI systems assisted with proof discovery, exposition, and an adversarial
review of generalized-inverse endpoint conventions.  Exact-rational checks
at selected parameter values and for small finite base laws were also used
during review as diagnostics.  No computational output is a premise of the
proof above; the quantified conclusion follows from the deductive argument.

\bibliographystyle{plain}
\bibliography{references}

\end{document}
